\section{Introduction}
\label{sec:intro}

La maintenance prédictive industrielle exige de plus en plus des modèles capables d'\emph{apprendre en
continu} sur l'équipement lui-même, sans retour vers un serveur : les conditions d'opération dérivent
(usure, changement de régime, de machine), et un modèle figé se dégrade. L'apprentissage continu, ou \gls{CL}, répond à ce besoin~\cite{DeLange2021Survey,Hurtado2023CLPdM}, mais son
déploiement sur un \gls{MCU} reste contraint : mémoire vive de l'ordre de la centaine de
kilo-octets, absence d'accélérateur neuronal, budget de latence strict.

Face à ces contraintes, la quantification INT8 est le réflexe le mieux établi de la littérature
embarquée~\cite{Ravaglia2021QLRCL,Capogrosso2023TinyML}. Elle est presque toujours justifiée par un
argument tripartite --- moins de mémoire, moins de temps, moins d'énergie --- rarement vérifié
\emph{sur la cible} pour chacun des trois termes séparément. C'est précisément ce que nous mesurons ici, et
les trois termes ne se comportent pas de la même manière.

Le présent travail se positionne sur un \textbf{triple gap} que peu de travaux comblent simultanément :
(Gap~1) validation sur des séries temporelles industrielles réelles ; (Gap~2) fonctionnement sous un budget
mémoire et de latence \emph{strict et mesuré} --- ici \SI{256}{\kilo\byte} de \gls{SRAM} et
\SI{100}{\milli\second} par cycle inférence + mise à jour ; (Gap~3) quantification INT8 compatible avec
l'entraînement incrémental. Nous instancions ces trois axes sur un classifieur
\gls{EWC}~\cite{Kirkpatrick2017EWC} porté en C sur une carte NUCLEO-F439ZI (Cortex-M4 \SI{180}{\mega\hertz},
\SI{256}{\kilo\byte} SRAM, sans \gls{NPU}), évalué sur deux jeux industriels aux scénarios continus contrastés.

\paragraph{Contributions.} (i)~Un protocole \emph{apparié} PC$\leftrightarrow$carte (même graine, même
ordre d'échantillons, même normalisation) établissant une parité FP32 \emph{exacte} en mode gelé
(\num{1.000}), quasi exacte en mode en ligne (de \num{0.9626} à \num{0.9887} sur les trois cellules),
pour une latence inférence +
apprentissage $\ll$~\SI{100}{\milli\second}. (ii)~La mise en évidence, \emph{sur carte réelle}, de
l'effondrement de la \gls{PTQ} naïve de la tête EWC binaire : le F1 tombe de $\approx\!\num{0.92}$ à \num{0.138}
et \num{0.1337} alors même que la \gls{RAM} des poids est bien divisée par~4. (iii)~L'isolation de la cause par
ablation \emph{émulée bit-exacte}, puis la \emph{récupération mesurée sur carte} avec un kernel à échelle
calibrée : le F1 revient à \num{0.9173} et \num{0.8995}, avec une parité gelée de \num{1.000} contre
l'émulateur et aucune erreur \gls{CRC}. (iv)~Deux axes complémentaires de quantification : le \emph{moment}, où
un entraînement \gls{QAT} n'ajoute pas de gain décisif face à une PTQ correctement calibrée, et la
\emph{profondeur}, où la frontière dépend de la métrique retenue --- ternaire en AUROC, INT4 en F1 mesuré
sur carte --- et où le gain mémoire n'existe qu'avec un véritable bit-packing. (v)~Un budget système complet mesuré sur la même carte, qui \emph{dissocie les trois
promesses de l'INT8} : la RAM des poids est bien divisée par~4, mais la RAM totale du système est
inchangée, la latence augmente d'environ la moitié et l'énergie ne bouge pas.

Ce dernier point est le fil conducteur de l'article. Sur cette cible, \textbf{l'INT8 se justifie par la
mémoire, ni par la vitesse ni par l'énergie} --- et encore, par une mémoire dont il faut préciser laquelle.
Nous distinguons systématiquement, dans le texte comme dans les légendes, ce qui est \emph{mesuré sur
carte} de ce qui est \emph{émulé bit-à-bit sur PC}, et laissons explicitement « à mesurer » toute cellule
non encore streamée, plutôt que de l'estimer.

\paragraph{Organisation.} La Section~\ref{sec:related} situe le travail, la Section~\ref{sec:method}
décrit le modèle, le protocole apparié et les schémas de quantification, la Section~\ref{sec:setup} le
dispositif. Les résultats se lisent en trois temps : parité, effondrement et récupération
(Section~\ref{sec:results}), moment et profondeur (Section~\ref{sec:quant-axes}), budget système
(Section~\ref{sec:system}). La Section~\ref{sec:discussion} discute les deux paradoxes et les limites.
